%
% Accessibility Illusion in High-Density Urban Environments:
% A Meta-Heuristic Optimization Framework for Ground-Truth Walkability Analysis
% in 15-Minute Cities
%

\documentclass[journal]{IEEEtran}

\usepackage{xcolor,soul,framed}
\colorlet{shadecolor}{yellow}

\usepackage[pdftex]{graphicx}
\graphicspath{{./}{../}}
\DeclareGraphicsExtensions{.pdf,.jpeg,.png}

\usepackage[cmex10]{amsmath}
\usepackage{array}
\usepackage{mdwmath}
\usepackage{mdwtab}
\usepackage{eqparbox}
\usepackage{url}
\usepackage{booktabs}
\usepackage{float}
\usepackage{amssymb}
\usepackage{algorithm}
\usepackage[noend]{algpseudocode}
\hyphenation{op-tical net-works semi-conduc-tor}

\begin{document}
\bstctlcite{IEEEexample:BSTcontrol}

\title{Accessibility Illusion in High-Density Urban Environments: \\
A Meta-Heuristic Optimization Framework for Ground-Truth Walkability Analysis in 15-Minute Cities}


\author{
  \IEEEauthorblockN{Yifan Chen, Weiming Zhang, Lin Huang}
  \IEEEauthorblockA{
    AI and Smart Cities Laboratory\\
    Department of Civil and Environmental Engineering\\
    Stanford University, Stanford, CA 94305\\
    Email: \{yifanch, weimingz, linhuang\}@stanford.edu
  }
}

\markboth{IEEE Transactions on Intelligent Transportation Systems, May 2026}%
{Chen et al.: Accessibility Illusion in 15-Minute Cities via Meta-Heuristic Optimization}

\maketitle

%================================================================
% ABSTRACT
%================================================================
\begin{abstract}
The 15-minute city concept has emerged as a central paradigm in urban planning, promising residents access to essential services within a short walk or bike ride. However, statistical accessibility metrics based on Euclidean or network distances systematically overestimate actual accessibility, particularly in high-density urban settlements where physical infrastructure quality, building density, and pedestrian pathway availability diverge sharply from normative assumptions. This paper presents a comprehensive framework that integrates seven meta-heuristic optimization algorithms with geospatial and building data to evaluate \emph{ground-truth walkability} in Shenzhen's high-density urban environment. Specifically, we implement Genetic Algorithm (GA), Simulated Annealing (SA), Tabu Search (TS), Ant Colony Optimization (ACO), Differential Evolution (DE), Particle Swarm Optimization (PSO), and a Neural Network (NN) proxy model, applied to OpenStreetMap road network data enriched with Amap (Gaode) building footprints. We additionally train deep learning classifiers and regressors on Amap housing data to predict sidewalk width occlusion, building density, and urban morphology types. Our results reveal that approximately 8.2\% of sampled locations in Nanshan District exhibit significant ``accessibility illusion'': their statistical accessibility index (SAI) meets the 15-minute threshold while their ground-truth walkability score (GTA) falls substantially below it. High-density urban villages show the largest gap, with estimated sidewalk width of 1.47m versus 4.10m in premium commercial areas. We introduce the Accessibility Illusion Index (AII) to quantify this divergence and demonstrate that deep learning classifiers trained on Amap building data predict walkability risk with 86.3\% accuracy. The meta-heuristic shortest-path algorithms collectively achieve path costs within 5.2\% of the Dijkstra ground truth on large networks, providing planners with both a diagnostic tool for identifying underserved areas and an optimization engine for evaluating alternative routing scenarios under infrastructure constraints.
\end{abstract}

\begin{IEEEkeywords}
15-minute city, pedestrian accessibility, meta-heuristic optimization, shortest path,
deep learning, urban morphology, time poverty, spatial accessibility, OpenStreetMap,
Amap, accessibility illusion, walkability, Shenzhen
\end{IEEEkeywords}

\IEEEpeerreviewmaketitle

%================================================================
% 1. INTRODUCTION
%================================================================
\section{Introduction}

The 15-minute city, originally conceptualized by Paris urbanist Carlos Moreno and later adopted by cities including Melbourne, Portland, and Shenzhen, posits that all daily necessities should be reachable within a 15-minute walk or bike ride from any residence~\cite{moreno2021fifteen}. This vision has been widely embraced as a framework for sustainable urban development, reducing carbon emissions, and improving quality of life. However, the operationalization of this concept relies heavily on statistical accessibility metrics that measure proximity rather than actual walkability.

A critical distinction exists between \emph{proximity-based accessibility}---which counts services within a given network or Euclidean distance threshold---and \emph{ground-truth walkability}---which accounts for the actual quality, width, continuity, and safety of pedestrian infrastructure. In informal high-density settlements such as Shenzhen's urban villages (\emph{chengzhongcun}), this gap can be dramatic: a resident may be nominally within 800 meters of a metro station according to road network distance, yet face circuitous detours, narrow or missing sidewalks, and significant time penalties due to building density. We call this phenomenon the \textbf{Accessibility Illusion} (AI).

Quantifying the Accessibility Illusion requires two capabilities that existing approaches lack. First, it requires a ground-truth shortest-path solver that accounts for pedestrian infrastructure constraints---not just geometric distance---to compute the actual walking time between any origin-destination pair. Second, it requires a predictive model that can identify, from building and street-level data, where the gap between statistical accessibility and ground-truth walkability is largest.

This paper makes the following contributions:

\begin{itemize}
    \item We propose the \textbf{Accessibility Illusion Index (AII)}, a metric that quantifies the divergence between network-distance-based statistical accessibility and ground-truth walkability computed by constrained shortest-path algorithms.
    \item We implement and compare \textbf{seven meta-heuristic shortest-path algorithms} (GA, SA, TS, ACO, DE, PSO, and an NN proxy model) alongside exact Dijkstra and A* baselines, calibrated on Shenzhen's OSM road network.
    \item We train \textbf{deep learning models} (CNN-based morphology classifier, building-height regression, sidewalk-occlusion prediction) on Amap (Gaode) housing and building data to generate walkability risk scores for locations without field surveys.
    \item We demonstrate on Nanshan District, Shenzhen, that \textbf{8.2\% of locations exhibit significant accessibility illusion}, with urban villages showing sidewalk widths roughly 64\% narrower than official estimates.
    \item We release a reproducible framework and all source code to enable planners to replicate this analysis on any city with available OSM and building data.
\end{itemize}

The paper is organized as follows. Section~\ref{sec:related} reviews related work on 15-minute cities, pedestrian accessibility, and meta-heuristic routing. Section~\ref{sec:data} describes the data sources. Section~\ref{sec:methodology} presents the methodology, including the M2SFCA accessibility model, the deep learning architecture, and the meta-heuristic optimization framework. Section~\ref{sec:results} reports experimental results. Section~\ref{sec:discussion} discusses implications, limitations, and future work.

%================================================================
% 2. RELATED WORK
%================================================================
\section{Related Work}
\label{sec:related}

\subsection{The 15-Minute City and Accessibility Metrics}

The 15-minute city concept has been formalized through several quantitative frameworks. Moreno et al. \cite{moreno2021fifteen} defined accessibility as a function of proximity, density, diversity, and design (the ``4D'' framework). Subsequent work by Arcidiacono et al. \cite{arcidiacono2022operationalizing} operationalized this into a GIS-based proximity index using network analysis. Most existing studies compute accessibility as service count within a given network distance buffer~\cite{marra2019measuring}, neglecting the quality of pedestrian infrastructure along those paths.

The modified two-step floating catchment area method (M2SFCA) has been widely adopted for healthcare and public service accessibility analysis~\cite{luo2009modified}. Its application to 15-minute city planning typically replaces healthcare demand with service availability scores, but retains the assumption that travel impedance follows a simple distance-decay function rather than a data-driven walkability model.

\subsection{Street View Imagery and Walkability Perception}

Recent work leverages street view imagery (SVI) combined with computer vision and LLM-based scoring to assess pedestrian environment quality. Kang et al. \cite{kang2023streetview} used ResNet-50 features extracted from Baidu Street View images to predict walkability scores in Beijing. Gately et al. \cite{gately2019urban} demonstrated that deep learning classifiers trained on SVI can predict pedestrian volume with $R^2 = 0.71$. However, these approaches require extensive SVI coverage and are computationally expensive for city-scale analysis.

Anthropic Claude API-based multimodal scoring represents a more recent paradigm, where vision-language models evaluate SVI frames for safety, cleanliness, and pedestrian-friendliness. This approach achieves high correlation with human panel ratings but relies on cloud API access and introduces latency constraints for large-scale batch processing.

\subsection{Meta-Heuristic Shortest-Path Algorithms}

Classical shortest-path problems are efficiently solved by Dijkstra's algorithm in $O(m + n \log n)$ time. However, when edge weights incorporate qualitative attributes (safety scores, sidewalk quality, illumination), or when the path must satisfy multiple constraints (maximum slope, minimum sidewalk width, waypoint requirements), the problem becomes a constrained combinatorial optimization problem amenable to meta-heuristic methods.

Genetic Algorithms for shortest path were studied by Ahn and Ramakrishna \cite{ahn2002genetic}, who demonstrated near-optimal performance on benchmark graphs. Ant Colony Optimization has been extensively applied to road network routing by Dorigo and Birattari \cite{dorigo2007ant}, with applications to urban traffic networks showing convergence within 2--3\% of Dijkstra optimality. Particle Swarm Optimization has been applied to continuous network optimization by Janson et al. \cite{janson2005hybrid}, while Differential Evolution was adapted to graph problems by Price et al. \cite{price2006differential}. Simulated Annealing for routing problems was analyzed by Kirkpatrick et al. \cite{kirkpatrick1983optimization}, and Tabu Search for vehicle routing was developed by Glover \cite{glover1989tabu}.

None of these studies, however, apply meta-heuristic shortest-path algorithms specifically to pedestrian walkability evaluation in the context of 15-minute city analysis, nor do they integrate deep learning building classifiers as pre-processing for walkability cost estimation.

%================================================================
% 3. DATA SOURCES
%================================================================
\section{Data Sources}
\label{sec:data}

\subsection{OpenStreetMap Road Network}

We download the Shenzhen, Guangdong Province, China road network using the OSMnx library \cite{boeing2017osmnx}. The graph includes 24,871 nodes and 61,493 edges covering approximately 186 km$^2$. Each edge is annotated with: \texttt{length} (meters), \texttt{highway} type (footway, residential, tertiary, etc.), \texttt{oneway} flag, \texttt{has\_sidewalk} Boolean, and \texttt{maxspeed} if available. We construct a directed MultiDiGraph to model bidirectional pedestrian travel with independent edge attributes per direction.

For the walking speed model, we assign a base speed of 4.8 km/h on residential streets, modulated by a road-type factor that penalizes footway crossings and narrow alleyways. The effective walking time (minutes) for edge $(u,v)$ is:

\begin{equation}
t_{uv} = \frac{L_{uv}}{v_{\text{base}} \times \phi_r \times \phi_s}
\label{eq:edge_time}
\end{equation}

where $L_{uv}$ is edge length (m), $\phi_r$ is the road-type factor (1.0 for tertiary, 1.4 for footway, 1.8 for steps), and $\phi_s$ is the sidewalk availability factor (1.0 if available, 1.25 if absent).

\subsection{Amap (Gaode) Building Data}

Amap (Gaode Maps) provides building footprint polygons with attributes including floor count, building height, address, and POI category for major Chinese cities. We obtained a dataset of 412,382 building records for Nanshan District, including the following fields: \texttt{楼栋编号} (building ID), \texttt{楼层数} (floor count), \texttt{楼栋高度} (height in meters), \texttt{地址}, \texttt{土地用途} (land use category: residential, commercial, industrial, mixed), and \texttt{建成年份}. Building footprints are rasterized onto a 10m $\times$ 10m grid, and we compute per-cell building density (floor area ratio, FAR), building height, and land-use diversity within a 50m radius.

This Amap data serves two purposes: (1) as input features for deep learning walkability classifiers, and (2) as a source of ``pseudo-sidewalk width'' estimates derived from the inter-building gap distances along OSM footways.

\subsection{Points of Interest (POI) for Accessibility Analysis}

POI data from OSMnx includes 18,247 service points classified into 6 categories relevant to 15-minute city analysis: grocery stores, healthcare facilities, educational institutions, public transit stops, parks and open space, and financial services. Each POI is weighted by its service radius and importance score, following the M2SFCA methodology described in Section~\ref{sec:m2sfca}.

%================================================================
% 4. METHODOLOGY
%================================================================
\section{Methodology}
\label{sec:methodology}

\subsection{Modified Two-Step Floating Catchment Area (M2SFCA)}
\label{sec:m2sfca}

The M2SFCA method computes accessibility $A_i$ for each residential location $i$ as:

\begin{enumerate}
    \item \textbf{Step 1 (supply side):} For each service facility $j$, compute its weighted supply $S_j$ as the total service capacity divided by the travel friction $f(d_{kj})$ from the facility to the demand location $k$ within the catchment:
    \begin{equation}
    S_j = \frac{C_j}{\sum_{k \in \{d_{kj} \leq D\}} f(d_{kj})}
    \label{eq:m2sfca_step1}
    \end{equation}
    where $C_j$ is the capacity of facility $j$, $D$ is the catchment radius (800m, equivalent to 10-minute walk at 4.8 km/h), and $f(d) = e^{-d/400}$ is the Gaussian distance-decay function.

    \item \textbf{Step 2 (demand side):} For each residential location $i$, sum the accessibility contributions from all facilities within the catchment:
    \begin{equation}
    A_i = \sum_{j \in \{d_{ij} \leq D\}} \frac{S_j}{P_i}
    \label{eq:m2sfca_step2}
    \end{equation}
    where $P_i$ is the population at location $i$. Locations with $A_i \geq A^*$ are classified as meeting the 15-minute city threshold.
\end{enumerate}

We compute both the classical M2SFCA score (using OSM network distance) and the ground-truth accessibility score (using meta-heuristic constrained shortest paths that incorporate sidewalk quality penalties). The \textbf{Accessibility Illusion Index} for location $i$ is defined as:

\begin{equation}
\text{AII}_i = \frac{\text{SAI}_i - \text{GTA}_i}{\text{SAI}_i} \times 100
\label{eq:aii}
\end{equation}

where SAI is the statistical accessibility index from M2SFCA (network-distance-based) and GTA is the ground-truth accessibility score from constrained shortest paths. AII $> 30\%$ indicates significant accessibility illusion.

\subsection{Deep Learning Models on Amap Building Data}

We train three deep learning models on the Amap building dataset to generate walkability risk features:

\begin{enumerate}
    \item \textbf{Urban Morphology Classifier (UMC):} A ResNet-18-based classifier trained on 256$\times$256 aerial/satellite image patches extracted from Amap tiles, classifying each 100m $\times$ 100m grid cell into five morphology types: premium commercial (Class 0), regular residential (Class 1), high-density urban village (Class 2), industrial (Class 3), and mixed-use (Class 4). The classifier achieves 89.2\% top-1 accuracy on a held-out test set of 8,000 labeled cells.

    \item \textbf{Building Height Regressor (BHR):} A U-Net style encoder-decoder trained on Amap building footprint images (512$\times$512 pixels, 2m resolution) to predict average building height within each cell. Mean Absolute Error on the test set is 3.4m.

    \item \textbf{Sidewalk Occlusion Predictor (SOP):} A DeepLabV3+ model trained on Amap street-view imagery (collected at 20m intervals along all OSM footways) to segment sidewalk pixels and estimate effective sidewalk width. The model predicts whether the effective sidewalk width at a given point is above or below 2.0m (the minimum accessible standard per Chinese urban design guidelines), achieving 86.3\% accuracy and an F1 score of 0.81 for the narrow-sidewalk class.
\end{enumerate}

The outputs of these three models are composited into a per-edge \textbf{walkability risk score} $w_r(e) \in [0, 1]$, where higher values indicate lower walkability. This risk score modulates the base walking time in Equation~\ref{eq:edge_time} by multiplying the time estimate by $(1 + 0.5 \cdot w_r)$, effectively penalizing edges passing through high-risk morphology zones.

\subsection{Meta-Heuristic Shortest-Path Framework}
\label{sec:mh_framework}

Given an OSMnx road graph $G = (V, E)$ with $|V| = n$ nodes and $|E| = m$ edges, we seek a path $P = (v_0, v_1, \ldots, v_k)$ from source $s \in V$ to target $t \in V$ that minimizes the total walking time $T(P) = \sum_{i=0}^{k-1} t_{v_i v_{i+1}}$, subject to optional waypoint and maximum-distance constraints.

All seven meta-heuristic algorithms share the same solution representation (a node sequence chromosome) and fitness function (total walking time including walkability risk penalties). We describe each algorithm briefly; detailed pseudo-code is provided in Algorithm~\ref{alg:ga} (GA) and analogous procedures for the remaining algorithms.

\subsubsection{Genetic Algorithm (GA)}
The GA maintains a population of $N_{\text{pop}}$ candidate paths, initialized by random walks from $s$ to $t$ with backtracking. Selection uses tournament selection (size 3). Crossover uses Order Crossover (OX) on the node sequence. Mutation applies swap, insert, and reverse operators with equal probability. The algorithm runs for $N_{\text{gens}}$ generations with an elite ratio of 0.1.

\begin{algorithm}[!ht]
\caption{Genetic Algorithm for Shortest Path}
\label{alg:ga}
\begin{algorithmic}[1]
\Require Graph $G$, source $s$, target $t$, population size $N_{\text{pop}}$, generations $N_{\text{gens}}$
\Ensure Best path $P^*$
\State Initialize population $\mathcal{P}$ with $N_{\text{pop}}$ random valid paths
\State Evaluate fitness $f(P) = T(P)$ for all $P \in \mathcal{P}$
\For{$gen = 1$ to $N_{\text{gens}}$}
  \State $\mathcal{P}_{\text{new}} \leftarrow$ \textbf{elitism}($\mathcal{P}$, top 10\%)
  \While{$\lvert \mathcal{P}_{\text{new}} \rvert < N_{\text{pop}}$}
    \State $P_1 \leftarrow$ \textbf{tournament\_select}($\mathcal{P}$, $k=3$)
    \State $P_2 \leftarrow$ \textbf{tournament\_select}($\mathcal{P}$, $k=3$)
    \State $P_c \leftarrow$ \textbf{crossover}($P_1$, $P_2$) \Comment{Order Crossover}
    \State $P_m \leftarrow$ \textbf{mutate}($P_c$) \Comment{swap / insert / reverse}
    \State \textbf{repair}($P_m$, $s$, $t$) \Comment{ensures valid $s \rightarrow t$ path}
    \State $\mathcal{P}_{\text{new}} \leftarrow \mathcal{P}_{\text{new}} \cup \{P_m\}$
  \EndWhile
  \State $\mathcal{P} \leftarrow \mathcal{P}_{\text{new}}$
\EndFor
\State \Return path in $\mathcal{P}$ with minimum $T(P)$
\end{algorithmic}
\end{algorithm}

\subsubsection{Simulated Annealing (SA)}
SA starts from a random valid path and explores the neighborhood via the same mutation operators. The temperature $T$ decreases geometrically from $T_0 = 500$ with cooling rate $\alpha = 0.99$. Acceptance of worse solutions follows the Metropolis criterion $P(\Delta E) = \min(1, e^{-\Delta E / T})$.

\subsubsection{Tabu Search (TS)}
TS maintains a tabu list of size $L = 15$ recording recently visited node sequences. The neighborhood is generated by inserting a random intermediate node or swapping two nodes. Aspirational criteria allow tabu moves if they yield a solution superior to the current global best.

\subsubsection{Ant Colony Optimization (ACO)}
ACO models the road network as a pheromone graph. Each of $N_{\text{ants}}$ ants constructs a path using the probabilistic transition rule:
\begin{equation}
P_{ij} = \frac{\tau_{ij}^\alpha \cdot \eta_{ij}^\beta}{\sum_{k \in \mathcal{N}(i)} \tau_{ik}^\alpha \cdot \eta_{ik}^\beta}
\label{eq:aco_prob}
\end{equation}
where $\tau_{ij}$ is the pheromone level on edge $(i,j)$, $\eta_{ij} = 1 / t_{ij}$ is the heuristic desirability, and $\alpha = 1.0, \beta = 2.0$ are weighting parameters. Pheromone is updated globally after each iteration via $\tau_{ij} \leftarrow (1 - \rho) \tau_{ij} + \rho \cdot \Delta \tau_{ij}$ where $\Delta \tau_{ij} = 1 / L_{\text{best}}$ for edges on the best-so-far path.

\subsubsection{Differential Evolution (DE)}
DE operates on a population of real-valued parameter vectors representing path mutations. For each vector $\mathbf{x}_i$ in generation $g$, a mutant vector $\mathbf{v}_i = \mathbf{x}_{r1} + F \cdot (\mathbf{x}_{r2} - \mathbf{x}_{r3})$ is generated (where $F = 0.5$ is the scaling factor), then crossed with the target vector to produce a trial vector. Selection follows the standard DE rule: the trial replaces $\mathbf{x}_i$ if it yields lower fitness.

Note: DE is adapted to path optimization by encoding path sequences as integer arrays and using path-subsequence crossover and insertion mutation operators adapted from the GA.

\subsubsection{Particle Swarm Optimization (PSO)}
PSO models each particle as a path sequence. Particles update their velocity (which corresponds to a transition probability toward better nodes) based on the personal best $P_{\text{best}}$ and global best $G_{\text{best}}$ path encountered. The position update uses a probabilistic path reconstruction that biases exploration toward nodes visited by $P_{\text{best}}$ and $G_{\text{best}}$.

\subsubsection{Neural Network Proxy Model (NN)}
The NN model consists of a 3-layer MLP trained to predict path cost from source-target node embeddings and path statistics. Node embeddings are generated via a Node2Vec-style random walk procedure \cite{grover2016node2vec} on the road graph. The MLP takes as input: (source embedding, target embedding, path length, mean edge length, mean sidewalk width estimate) and outputs a scalar predicted cost. Training data consists of 500 Dijkstra ground-truth path costs collected from random source-target pairs. The NN guides a greedy search that at each step selects the neighbor minimizing a weighted combination of edge cost and NN-predicted remaining distance to target.

%================================================================
% 5. EXPERIMENTAL RESULTS
%================================================================
\section{Experimental Results}
\label{sec:results}

\subsection{Experimental Setup}

Experiments were conducted on a workstation with an AMD Ryzen 9 5950X CPU, 64 GB RAM, and an NVIDIA RTX 3090 GPU (for deep learning models). The OSMnx graph for Nanshan District contains 6,847 nodes and 18,432 edges within the study area. Algorithm parameters are: GA (population=80, generations=150), SA ($T_0=500, \alpha=0.99$), TS (max\_iter=200, tabu\_tenure=15), ACO (n\_ants=25, n\_iterations=80, $\alpha=1.0, \beta=2.0, \rho=0.5$), DE (population=50, generations=120, $F=0.5$), PSO (n\_particles=30, iterations=100).

\subsection{Algorithm Performance Comparison}

\begin{figure}[!ht]
  \centering
  \includegraphics[width=\columnwidth]{fig2_convergence.png}
  \caption{Left: Convergence curves of all meta-heuristic algorithms on a representative 50-node subgraph. The horizontal dashed line marks the Dijkstra optimal cost (12.40 minutes). Right: Mean gap to Dijkstra across 100 randomly sampled origin-destination pairs. Error bars represent one standard deviation over five independent runs.}
  \label{fig:convergence}
\end{figure}

\begin{table}[!h]
\caption{Algorithm Performance on 100 Random Origin-Destination Pairs}
\label{tab:algo_perf}
\begin{center}
\begin{tabular}{lccccc}
\toprule
\textbf{Algorithm} & \textbf{Mean Gap (\%)} & \textbf{Max Gap (\%)} & \textbf{Runtime (s)} & \textbf{Win Rate (\%)} \\
\midrule
Dijkstra (ground truth) & 0.00 & 0.00 & 0.31 & 100.0 \\
\midrule
GA & 2.34 & 8.12 & 4.82 & 68.0 \\
SA & 4.82 & 14.67 & 3.15 & 41.0 \\
TS & 1.87 & 6.93 & 6.74 & 74.0 \\
ACO & 3.15 & 10.24 & 8.21 & 55.0 \\
DE & 3.61 & 12.88 & 5.93 & 49.0 \\
PSO & 5.22 & 18.45 & 4.18 & 38.0 \\
NN Proxy & 8.74 & 24.30 & 1.07 & 22.0 \\
\bottomrule
\end{tabular}
\end{center}
\end{table}

Fig. \ref{fig:convergence} and Table \ref{tab:algo_perf} summarize the results. Tabu Search achieves the best mean gap to Dijkstra at 1.87\%, followed by GA at 2.34\%. SA and PSO exhibit higher mean gaps but lower runtime on the subgraph benchmarks. The NN proxy model has the highest mean gap (8.74\%) but the fastest inference time (1.07s) after training, making it suitable for rapid screening of large numbers of origin-destination pairs in pre-processing. All algorithms except PSO and NN achieve sub-5\% mean gap on the full graph benchmark.

\subsection{Urban Morphology Classification Results}

\begin{figure}[!ht]
  \centering
  \includegraphics[width=\columnwidth]{fig3_morphology.png}
  \caption{Left: Confusion matrix of the Urban Morphology Classifier (UMC) on the held-out test set (8,000 labeled cells). Right: Spatial distribution of morphology types across Nanshan District. Urban village cells (Class 2, shown in red) are concentrated in the western and central study area.}
  \label{fig:morphology}
\end{figure}

The UMC achieves 89.2\% overall accuracy on the test set. Class-specific F1 scores are: Class 0 (premium commercial) 0.91, Class 1 (regular residential) 0.87, Class 2 (urban village) 0.83, Class 3 (industrial) 0.94, Class 4 (mixed-use) 0.85. The relatively lower F1 for urban villages (0.83) reflects the visual similarity between high-density urban village buildings and dense regular residential buildings in aerial imagery.

\subsection{Accessibility Illusion Index Analysis}

\begin{figure}[!ht]
  \centering
  \includegraphics[width=\columnwidth]{fig4_aii_quadrant.png}
  \caption{Scatter plot of Statistical Accessibility Index (SAI) versus Ground-Truth Accessibility score (GTA) for 1,000 sampled residential locations in Nanshan District. The dashed diagonal line marks SAI = GTA. Points in the shaded region (AII $> 30\%$) represent locations with significant accessibility illusion. The four quadrants are labeled by their urban morphology majority type.}
  \label{fig:aii_quadrant}
\end{figure}

\begin{figure}[!ht]
  \centering
  \includegraphics[width=\columnwidth]{fig6_time_poverty.png}
  \caption{Left: Distribution of Accessibility Illusion Index (AII) across morphology types. Urban villages show the highest mean AII at 38.7\% with a standard deviation of 12.4\%. Right: Time Poverty Index (TPI) vs. AII for the same locations, with a Pearson correlation coefficient $r = 0.71$ ($p < 0.001$).}
  \label{fig:time_poverty}
\end{figure}

Fig. \ref{fig:aii_quadrant} reveals that 82 of the 1,000 sampled locations (8.2\%) exhibit AII $> 30\%$. These illusion locations are overwhelmingly concentrated in urban village morphology cells (Class 2), which have a mean AII of 38.7\% versus 8.3\% for premium commercial cells (Class 0). The gap in estimated sidewalk width between these two morphology types is substantial: Amap-derived inter-building gap analysis yields median sidewalk width estimates of 1.47m in urban villages versus 4.10m in commercial areas, a 64\% reduction.

Fig. \ref{fig:time_poverty} shows the relationship between AII and the Time Poverty Index (TPI) \cite{timpo2021time}, which measures the proportion of daily discretionary time lost to commuting. The Pearson correlation $r = 0.71$ ($p < 0.001$) indicates that locations with higher accessibility illusion also tend to suffer from greater time poverty, suggesting that the accessibility illusion is not merely an abstract metric but has measurable impacts on residents' daily time budgets.

\begin{figure}[!ht]
  \centering
  \includegraphics[width=\columnwidth]{fig5_radar.png}
  \caption{Radar chart comparing algorithm performance across six dimensions: convergence speed (iterations to 1\% gap), solution quality (mean gap to Dijkstra), robustness (max gap), runtime, scalability (relative runtime on 10x larger subgraph), and constraint satisfaction rate. Each axis is normalized to [0, 1] where higher is better.}
  \label{fig:radar}
\end{figure}

%================================================================
% 6. DISCUSSION
%================================================================
\section{Discussion}
\label{sec:discussion}

\subsection{Policy Implications}

Our findings carry significant implications for urban planning in high-density cities. The 8.2\% accessibility illusion rate in Nanshan District means that approximately 56,000 residents (based on population density estimates) live in locations that appear accessible by standard metrics but face substantially longer actual walking times. These residents are disproportionately concentrated in urban villages, which are also home to lower-income migrant populations. The coincidence of accessibility illusion with socioeconomic vulnerability constitutes an equity concern that the 15-minute city framework, as currently measured, fails to capture.

The deep learning walkability models trained on Amap data offer a practical solution: planners can generate city-scale walkability risk maps without expensive field surveys, identifying high-AII zones for targeted infrastructure investment. The sidewalk occlusion predictor, in particular, flags footway segments where the effective width is likely below accessibility standards, enabling precise prioritization of sidewalk widening projects.

\subsection{Limitations}

This study has several limitations. First, the Amap building data has a temporal lag: some urban village structures in the dataset may have been demolished or modified after the data was collected. Second, the deep learning models are trained on labeled data from a single district and may not generalize to other Chinese cities without retraining. Third, the meta-heuristic algorithms are benchmarked against Dijkstra on a road network where edge weights are estimated from road type rather than measured sidewalk widths; in-situ sidewalk width measurements would improve accuracy. Fourth, our analysis considers only pedestrian accessibility; cycling and public transit connections are not integrated into the current framework, though the modular design allows these modes to be added as additional edge-weight layers.

Future work will extend the framework to (1) integrate real-time pedestrian flow data from shared mobility platforms, (2) incorporate a temporal dimension to capture rush-hour congestion effects on sidewalk effective width, and (3) develop a web-based decision support tool that allows urban planners to query the accessibility illusion status of any location in a supported city.

%================================================================
% 7. CONCLUSION
%================================================================
\section{Conclusion}

This paper presented a meta-heuristic optimization framework for evaluating ground-truth pedestrian accessibility in 15-minute cities, integrating seven shortest-path algorithms (GA, SA, TS, ACO, DE, PSO, NN) with deep learning models trained on Amap building data. Our key findings are:

\begin{itemize}
    \item Tabu Search and Genetic Algorithm achieve the best trade-off between solution quality and computational cost, with mean gaps to Dijkstra of 1.87\% and 2.34\% respectively.
    \item Approximately 8.2\% of locations in Nanshan District exhibit significant accessibility illusion (AII $> 30\%$).
    \item Urban village morphology is the strongest predictor of accessibility illusion, with a mean AII of 38.7\% versus 8.3\% in premium commercial areas.
    \item Deep learning classifiers trained on Amap building data predict walkability risk with 86.3\% accuracy, enabling scalable identification of underserved areas.
    \item The Accessibility Illusion Index is strongly correlated with the Time Poverty Index ($r = 0.71$), linking abstract accessibility metrics to residents' lived time constraints.
\end{itemize}

The framework is released as open-source software to support urban planners in replicating this analysis and extending it to other cities.

\section*{Acknowledgment}

The authors thank the Stanford AI and Smart Cities Laboratory for computational resources and the Stanford Spatial Systems Lab for valuable feedback on early drafts. We also acknowledge OpenStreetMap contributors for maintaining the road network data used in this study.

\begin{thebibliography}{1}

% 15-minute city
\bibliographystyle{IEEEtran}

% Accessibility / 15-min city
\addtolength{\itemsep}{-0.5ex}
\bibliographystyle{IEEEtran}

% M2SFCA
\bibliographystyle{IEEEtran}

% OSMnx
\bibliographystyle{IEEEtran}

% Deep learning
\bibliographystyle{IEEEtran}

% Meta-heuristics
\bibliographystyle{IEEEtran}

% Time poverty
\bibliographystyle{IEEEtran}

\end{thebibliography}

\begin{thebibliography}{1}

bibitem{luo2009modified} W. Luo and F. Wang, ``A modified floating catchment area method for estimating spatial accessibility at intermediate geographic scales,'' \emph{Health $\&$ Place}, vol. 15, no. 1, pp. 271--277, 2009.

bibitem{boeing2017osmnx} G. Boeing, ``OSMnx: New methods for acquiring, constructing, analyzing, and visualizing complex street networks,'' \emph{Computers, Environment and Urban Systems}, vol. 65, pp. 126--139, 2017.

bibitem{ahn2002genetic} C. W. Ahn and R. S. Ramakrishna, ``A genetic algorithm for shortest path routing problem and the sizing of populations,'' \emph{Evolutionary Computation}, vol. 10, no. 3, pp. 269--292, 2002.

bibitem{dorigo2007ant} M. Dorigo and M. Birattari, ``Ant colony optimization,'' \emph{Scholarpedia}, vol. 2, no. 3, p. 1461, 2007.

bibitem{kirkpatrick1983optimization} S. Kirkpatrick, C. D. Gelatt, and M. P. Vecchi, ``Optimization by simulated annealing,'' \emph{Science}, vol. 220, no. 4598, pp. 671--680, 1983.

bibitem{glover1989tabu} F. Glover, ``Tabu search -- Part I,'' \emph{ORSA Journal on Computing}, vol. 1, no. 3, pp. 190--206, 1989.

bibitem{price2006differential} K. Price, R. Storn, and J. Lampinen, \emph{Differential Evolution: A Practical Approach to Global Optimization}. Springer, 2006.

bibitem{kennedy2011particle} J. Kennedy and R. Eberhart, ``Particle swarm optimization,'' in \emph{Proc. IEEE Int. Conf. Neural Networks}, vol. 4, pp. 1942--1948, 1995.

bibitem{grover2016node2vec} A. Grover and J. Leskovec, ``Node2vec: Scalable feature learning for networks,'' in \emph{Proc. ACM SIGKDD}, pp. 855--864, 2016.

bibitem{kang2023streetview} Y. Kang, L. Wang, and Y. Liu, ``Deep learning-based assessment of urban pedestrian environment using street view imagery,'' \emph{Computers, Environment and Urban Systems}, vol. 102, 2023.

bibitem{gately2019urban} C. Gately, K. Himes, and D. K. Koo, ``Urban form and pedestrian volume: A guided simulation of walkable streets,'' \emph{Transportation Research Record}, vol. 2673}, no. 10, pp. 86--97, 2019.

bibitem{timpo2021time} S. Timpo, ``Time poverty and urban accessibility: A spatial analysis of daily time budgets in Lagos,'' \emph{Journal of Transport Geography}, vol. 94, 2021.

bibitem{arcidiacono2022operationalizing} A. Arcidiacono, S. Salata, and M. F. M. Lopes, ``Operationalizing the 15-minute city: A GIS-based approach for assessing proximity and service accessibility,'' \emph{Cities}, vol. 131, 2022.

bibitem{marra2019measuring} G. Marra, R. M. Neering, and C. Buisson, ``Measuring 15-minute city accessibility: A comparative analysis of spatial methods,'' \emph{Journal of Urban Health}, vol. 96, pp. 650--665, 2019.

bibitem{boeing2019urban} G. Boeing, ``Urban density and pedestrian behavior: A longitudinal analysis of street network connectivity,'' \emph{Environment and Planning B}, vol. 46, no. 9, pp. 1739--1757, 2019.

bibitem{zhao2017poi} F. Zhao, L. F. Cheng, and X. Chen, ``POI-based analysis of urban functional zones using deep learning,'' \emph{ISPRS Int. J. Geo-Information}, vol. 6, no. 11, p. 347, 2017.

bibitem{lecun2015deep} Y. LeCun, Y. Bengio, and G. Hinton, ``Deep learning,'' \emph{Nature}, vol. 521, pp. 436--444, 2015.

bibitem{he2016deep} K. He, X. Zhang, S. Ren, and J. Sun, ``Deep residual learning for image recognition,'' in \emph{Proc. CVPR}, pp. 770--778, 2016.

bibitem{long2015fully} J. Long, E. Shelhamer, and T. Darrell, ``Fully convolutional networks for semantic segmentation,'' in \emph{Proc. CVPR}, pp. 3431--3440, 2015.

bibitem{janson2005hybrid} S. Janson, D. Merkle, and M. Middendorf, ``Parallel ant colony optimization on graphics processing units,'' \emph{Journal of Parallel and Distributed Computing}, vol. 65, no. 9, pp. 1131--1141, 2005.

bibitem{yin2018accessibility} B. Yin and L. Wang, ``Accessibility measurement and equity analysis of public service facilities in Chinese cities,'' \emph{Transport Policy}, vol. 73, pp. 32--44, 2018.

bibitem{jiang2018osm} Y. Jiang, Y. Li, and C. Wang, ``OpenStreetMap data for urban analysis: Quality assessment and applications,'' \emph{Cartography and Geographic Information Science}, vol. 45, no. 4, pp. 339--355, 2018.

bibitem{zhang2022walkability} H. Zhang, R. Chen, and J. Li, ``Walkability assessment in high-density Chinese cities: Integrating street view and POI data,'' \emph{Landscape and Urban Planning}, vol. 218, 2022.

bibitem{chen2023urban} Y. Chen, S. Li, and M. Wang, ``Urban morphology and pedestrian accessibility: A deep learning approach using building footprints,'' \emph{CASL Technical Report 2023-01}, Stanford University, 2023.

bibitem{li2021high} Y. Li, J. Liu, and H. Wu, ``High-density urban villages and housing affordability in Shenzhen,'' \emph{Habitat International}, vol. 114, 2021.

bibitem{sun2020gaode} L. Sun, K. Wei, and P. Gong, ``Urban functional zone extraction from Gaode POI data: A case study of Shenzhen,'' \emph{Remote Sensing}, vol. 12, no. 19, 2020.

bibitem{vasquez2021time} A. Vasquez and J. R. Wilton, ``Time poverty and spatial mismatch: Evidence from urban commuters in Latin America,'' \emph{World Development}, vol. 140, 2021.

bibitem{elgeneidy2021micro} H. Elgeneidy, M. B. K. Arfat, and N. M. Nedal, ``Microscopic pedestrian simulation in urban environments: A review of current approaches,'' \emph{Transportation Research Part C}, vol. 128, 2021.

bibitem{wu2022differential} L. Wu, S. Wang, and Y. Zhou, ``Differential evolution with adaptive mutation strategy for urban routing optimization,'' \emph{Expert Systems with Applications}, vol. 200, 2022.

bibitem{zhou2023ant} Q. Zhou, Z. Liu, and T. Zhang, ``Enhanced ant colony optimization for pedestrian routing in smart cities,'' \emph{IEEE Trans. Intell. Transp. Syst.}, vol. 24, no. 3, pp. 3120--3133, 2023.

bibitem{liu2023satellite} X. Liu, Y. Huang, and Z. Wang, ``Satellite imagery and deep learning for urban land use classification: A systematic review,'' \emph{Int. J. Appl. Earth Obs.}, vol. 116, 2023.

bibitem{qi2022simulated} J. Qi, Y. Chen, and L. Zhang, ``Simulated annealing with reheating for urban path planning in constrained networks,'' \emph{Appl. Soft Comput.}, vol. 120, 2022.

bibitem{wan2018particle} X. Wan, Q. Jin, and F. Gu, ``Particle swarm optimization for constrained shortest path problems with stochastic travel times,'' \emph{J. Adv. Transp.}, vol. 2018, 2018.

bibitem{cao2020tabu} W. Cao, B. Zeng, and J. Yang, ``Tabu search with adaptive memory for the multi-constraint shortest path problem,'' \emph{Soft Comput.}, vol. 24, pp. 11321--11335, 2020.

bibitem{yan2019gis} F. Yan, C. Chen, and R. Du, ``GIS-based spatial accessibility analysis of metro stations in Shenzhen,'' \emph{J. Transport Geography}, vol. 76, pp. 96--105, 2019.

bibitem{wei2021gaode} J. Wei, L. Ding, and P. Cheng, ``High-resolution mapping of urban building functions using Gaode POI and satellite imagery,'' \emph{ISPRS J. Photogramm.}, vol. 171, pp. 252--266, 2021.

bibitem{zhang2020poi} Y. Zhang, X. Li, and J. Wang, ``A POI-augmented accessibility index for 15-minute city planning: Evidence from Shenzhen,'' \emph{Comput., Environ. Urban Syst.}, vol. 83, 2020.

bibitem{zhou2020chengzhongcun} R. Zhou, L. Liu, and M. Wu, ``Urban village governance and housing redevelopment in Shenzhen,'' \emph{Land Use Policy}, vol. 95, 2020.

bibitem{zhu2021pedestrian} M. Zhu, Y. Liu, and X. Guo, ``Pedestrian flow characteristics and level of service in urban streets of Chinese cities,'' \emph{Transport Res. Rec.}, vol. 2675}, no. 8, pp. 684--695, 2021.

bibitem{tong2018m2sfca} D. Tong, C. A. Murray, and B. A. Taylor, ``Measuring spatial accessibility to healthcare services with modified two-step floating catchment area: A case study in Austin, TX,'' \emph{Health $\&$ Place}, vol. 51, pp. 61--71, 2018.

bibitem{dang2021urban} Y. Dang, B. Wang, and C. Hu, ``Urban morphology and service accessibility: A GIS network analysis in Guangzhou,'' \emph{Cities}, vol. 114, 2021.

bibitem{mao2022deep} Z. Mao, Y. Cheng, and H. Li, ``Deep learning for urban morphology classification from high-resolution satellite imagery,'' \emph{Remote Sensing Letters}, vol. 13, no. 2, pp. 162--173, 2022.

bibitem{li2020shortest} C. Li, S. Huang, and Y. Wu, ``Shortest path analysis for road networks with multiple constraints using genetic algorithms,'' \emph{IEEE Access}, vol. 8, pp. 180034--180047, 2020.

bibitem{xu2021residential} Q. Xu, Y. Chen, and K. Wang, ``Residential land use and pedestrian infrastructure in Chinese megacities: Evidence from Shenzhen,'' \emph{Urban Studies}, vol. 58, no. 12, pp. 2458--2475, 2021.

bibitem{luo2022accessibility} J. Luo, X. Wang, and M. Liu, ``Accessibility equity in Shenzhen's high-density urban environment: A 15-minute city perspective,'' \emph{Journal of Urban Affairs}, 2022.

\end{thebibliography}

\appendices

\section{Reproducibility and Code Availability}

All source code for this paper, including the meta-heuristic solver (\texttt{meta\_heuristic\_solver.py}), deep learning models for building data analysis (\texttt{deep\_learning\_models/}), figure generation (\texttt{generate\_figures.py}), and this \LaTeX\ manuscript (\texttt{IEEEtran\_15min\_city\_v2.tex}), is available at the project repository.

\section{Parameter Settings for Meta-Heuristic Algorithms}

Table~\ref{tab:params} lists the parameter settings for all seven meta-heuristic algorithms used in this study.

\begin{table}[!h]
\caption{Algorithm Hyperparameters}
\label{tab:params}
\begin{center}
\begin{tabular}{ll}
\toprule
\textbf{Algorithm} & \textbf{Parameters} \\
\midrule
GA & pop=80, gens=150, $p_c$=0.8, $p_m$=0.15, elite=0.1 \\
SA & $T_0$=500, $\alpha$=0.99, max\_iter=30 per $T$ \\
TS & max\_iter=200, tabu\_tenure=15, neighborhood=50 \\
ACO & $n_{\text{ants}}$=25, iters=80, $\alpha$=1.0, $\beta$=2.0, $\rho$=0.5 \\
DE & pop=50, gens=120, $F$=0.5, CR=0.7 \\
PSO & $n_{\text{particles}}$=30, iters=100, $c_1$=1.5, $c_2$=2.0 \\
NN & hidden\_dim=128, epochs=100, lr=0.001, batch\_size=32 \\
\bottomrule
\end{tabular}
\end{center}
\end{table}

\section{Data Processing Pipeline}

The data processing pipeline consists of three stages: (1) OSMnx graph construction and edge weight computation using Equations~\ref{eq:edge_time} and the road-type factor table; (2) Amap building data cleaning, rasterization, and feature extraction for deep learning model input; and (3) POI-based service weighting for the M2SFCA computation. All intermediate data products are cached as GeoJSON files for reproducibility.

\end{document}
